\section{Data Transformation Pipeline}

\subsection{Band Extraction from S3 COGs}

When a query is issued for a specific spatiotemporal window, the system first executes the PostGIS spatiotemporal query to retrieve matching scene URLs. For each scene, the Red (Band 1), NIR (Band 2), and SWIR1 (Band 3) band arrays are read directly from S3 COGs using \texttt{rasterio} with an AWS Session, streaming only the required bands without downloading complete files. A configurable \texttt{pixel\_stride} parameter (default 10) enables subsampling for scalability evaluations at reduced resolution. The \texttt{rasterio} dataset mask is applied to set no-data pixels to NaN. The strided affine transform is computed to preserve accurate pixel centroid coordinates for subsequent geospatial operations. Scene date is parsed from the file path or \texttt{tile\_id} using regular expression matching across multiple HLS and custom filename patterns.

\subsection{Forest Mask Application}

Following band extraction, the ESA WorldCover forest mask is applied at the NumPy array level before PySpark ingestion. Using \texttt{rasterio}'s \texttt{reproject} function, the WorldCover GeoTIFF is reprojected onto the subsampled HLS pixel grid, and pixels with values not equal to class 10 (tree cover) are set to NaN. This filtering reduces the total pixel count substantially: in typical Boulder County scenes, forested pixels represent approximately 30--70\% of the study area extent. Performing the masking at the NumPy level before Spark ingestion minimizes Spark shuffle overhead associated with filtering null rows.

\subsection{Vegetation Index Computation in PySpark}

NDVI and NDMI are computed within PySpark using column expressions on the pixel-level DataFrame:

\begin{equation}
\text{NDVI} = \frac{\text{NIR} - \text{Red}}{\text{NIR} + \text{Red}}
\end{equation}

\begin{equation}
\text{NDMI} = \frac{\text{NIR} - \text{SWIR1}}{\text{NIR} + \text{SWIR1}}
\end{equation}

Both formulas are guarded against division-by-zero by a Spark \texttt{WHEN} clause that returns \texttt{NULL} when the denominator equals zero. NDVI values approaching $+1$ indicate dense healthy vegetation; NDMI values approaching $+1$ indicate high canopy moisture content. These two complementary indices capture both photosynthetic activity and water status, providing richer characterization of vegetation stress than either alone, as motivated by the spectral analysis literature~\cite{carletti2025,konig2023}.

\subsection{Temporal Harmonization and Monthly Compositing}

Irregular multi-sensor observations are standardized into monthly composites using a Spark \texttt{groupBy} aggregation on stable pixel IDs (\texttt{true\_pixel\_id}), year, and month. Mean NDVI (\texttt{ndvi\_mean}) and mean NDMI (\texttt{ndmi\_mean}) are computed per pixel-month. Pixel-months with all-NaN values are dropped. The monthly aggregation serves as a noise-reduction step: the mean over multiple observations within a month smooths sensor-specific artifacts and minor residual cloud contamination. The resulting regular time series is the direct input to the Z-score anomaly detection component.